\section*{Chapitre \ref{ch:b1:plate-kinematics} --- Cinématique des plaques sur une sphère}

\begin{solution}{exo:b1:plate-kinematics:1}
(a) $\qty{0.210}{\degree\per\Ma} = 0.210\times\pi/180 = \qty{3.665e-3}{rad}$
par million d'années, c'est-à-dire \qty{3.665e-9}{rad.an^{-1}}~; avec
$\qty{1}{\year} = \qty{3.156e7}{s}$, \qty{1.161e-16}{rad.s^{-1}}. (b) $360/0.210 =
\qty{1714}{\Ma}$, plus d'un tiers de l'âge de la Terre. (c) $\omega R =
\num{3.665e-9}\times\qty{6.371e9}{mm} = \qty{23.35}{mm.\year^{-1}}$.
\end{solution}

\begin{solution}{exo:b1:plate-kinematics:2}
$v = \omega R\sin\Delta = 23.35\sin\Delta$~: \qty{11.7}{mm.an^{-1}} à \ang{30},
\qty{20.2}{mm.an^{-1}} à \ang{60} et \qty{23.35}{mm.an^{-1}} à \ang{90}. L'expansion
est la plus rapide là où la limite est la plus éloignée du pôle, jusqu'à
\ang{90}~: entre l'Eurasie et l'Amérique du Nord, à l'extrémité sud de la
limite, près des Açores, à \ang{78} du pôle.
\end{solution}

\begin{solution}{exo:b1:plate-kinematics:3}
Le sous-chron a duré $1.072 - 0.988 = \qty{0.084}{\Ma}$, et une demi-vitesse en
\unit{mm.an^{-1}} est aussi une vitesse en \unit{km\per\Ma}. (a) $50\times0.084 =
\qty{4.2}{km}$ à \qty{50}{mm.an^{-1}}~; $10\times0.084 = \qty{0.84}{km}$ à
\qty{10}{mm.an^{-1}}. (b) La bande chevauche l'axe~:
$2\times10\times0.781 = \qty{15.6}{km}$.
\end{solution}

\begin{solution}{exo:b1:plate-kinematics:4}
Les demi-vitesses sont $66/2.581 = \qty{25.6}{mm.\year^{-1}}$ vers l'ouest et
$64/2.581 = \qty{24.8}{mm.\year^{-1}}$ vers l'est~; la vitesse totale est
$130/2.581 = \qty{50.4}{mm.\year^{-1}}$.
L'expansion est presque symétrique~: le flanc ouest a crû d'environ
\qty{3}{\percent} plus vite.
\end{solution}

\begin{solution}{exo:b1:plate-kinematics:5}
$\cos\Delta = \sin\ang{61.8}\sin\ang{85} + \cos\ang{61.8}\cos\ang{85}
\cos(\ang{0} - \ang{139.6}) = 0.8813\times0.9962 + 0.4726\times0.08716\times
(-0.7615) = 0.87795 - 0.03137 = 0.84658$, donc $\Delta = \ang{32.16}$ et
$v = 23.35\sin\ang{32.16} = \qty{12.4}{mm.\year^{-1}}$. C'est un peu au-dessus des
\qtyrange{11}{12}{mm.an^{-1}} mesurés, mais le \omterm{def:b1:plate-kinematics:euler-pole}{vecteur d'Euler} est un unique meilleur
ajustement de 453 vitesses, dont les vitesses s'écartent de \qty{0.7}{mm.an^{-1}} en
moyenne, et les vitesses individuelles près de la lente dorsale de Gakkel sont parmi les moins
précises.
\end{solution}

\begin{solution}{exo:b1:plate-kinematics:6}
(a) $\cos\Delta = \sin\ang{61.8}\sin\ang{39.44} + \cos\ang{61.8}\cos\ang{39.44}
\cos(\ang{169.43}) = 0.20111$~: $\Delta = \ang{78.40}$. (b) $v = 23.35
\sin\ang{78.40} = \qty{22.87}{mm.\year^{-1}}$, les \qty{23}{mm.an^{-1}} mesurés
juste au nord du \omterm{def:b1:plate-kinematics:triple-junction}{point triple}. (c) Avec la formule de la méthode, le numérateur
est $\sin\ang{169.43}\cos\ang{61.8} = 0.08668$ et le dénominateur
$\cos\ang{39.44}\sin\ang{61.8} - \sin\ang{39.44}\cos\ang{61.8}\cos\ang{169.43}
= 0.97573$~: le pôle est vu à $\beta = \ang{5.08}$, presque plein nord, et
l'Eurasie se déplace vers $\ang{95.1}$. (d) Les transformantes sont parallèles au
mouvement relatif~: elles sont orientées vers \ang{095}, presque est--ouest.
\end{solution}

\begin{solution}{exo:b1:plate-kinematics:7}
(a) $\vec\omega = \omega(\cos\varphi_p\cos\lambda_p, \cos\varphi_p\sin\lambda_p,
\sin\varphi_p)$~: $\vec\omega_{\mathrm{EU\,NA}} = (-0.0756, 0.0643, 0.1851)$ et
$\vec\omega_{\mathrm{NB\,NA}} = (0.0333, 0.0282, 0.2289)$, en
\unit{\degree\per\Ma}. (b) $\vec\omega_{\mathrm{NB\,EU}} =
\vec\omega_{\mathrm{NB\,NA}} - \vec\omega_{\mathrm{EU\,NA}} = (0.1089, -0.0361,
0.0438)$, de longueur
\[
  \sqrt{0.011859 + 0.001303 + 0.001918} = \qty{0.1228}{\degree\per\Ma}.
\]
Son pôle est à la \omterm{def:b1:spherical-astronomy:geographic}{latitude}
$\arcsin(0.0438/0.1228) = \ang{20.9}$\,N et à la \omterm{def:b1:spherical-astronomy:geographic}{longitude}
$\arctan(-0.0361/0.1089) = \ang{-18.3}$, c'est-à-dire \ang{18.3}\,O (la composante
$x$ étant positive, on n'ajoute pas \ang{180}). (c) La limite entre
la Nubie et l'Eurasie offre peu de mesures directes~: entre les Açores et
Gibraltar, seulement les \omterm{def:b1:spherical-astronomy:horizontal}{azimuts} de la faille Gloria. Les deux mouvements d'expansion
à travers l'Atlantique sont mesurés par des centaines de vitesses, et leur différence
hérite de leur précision.
\end{solution}

\begin{solution}{exo:b1:plate-kinematics:8}
(a) Autour de $\vec e_z$ de \ang{90}~: $\cos\ang{90} = 0$, $\vec e_z\cdot\vec e_x = 0$
et $\vec e_z\times\vec e_x = \vec e_y$, donc $\vec e_x$ va en $\vec e_y$. Puis
autour de $\vec e_x$ de \ang{90}~: $\vec e_x\cdot\vec e_y = 0$ et
$\vec e_x\times\vec e_y = \vec e_z$, donc $\vec e_y$ va en $\vec e_z$~: le point
finit au pôle nord. (b) D'abord autour de $\vec e_x$~: $\vec e_x$ est sur l'axe,
$\vec e_x\times\vec e_x = \vec 0$ et $(1 - \cos\ang{90})(\vec e_x\cdot\vec e_x)
\vec e_x = \vec e_x$, donc il ne bouge pas~; puis autour de $\vec e_z$ il va en
$\vec e_y$, à \ang{0}\,N, \ang{90}\,E. (c) Les résultats diffèrent~: les \omterm{def:b1:plate-kinematics:finite-rotation}{rotations
finies} ne commutent pas. Des reconstructions successives doivent se composer en
multipliant les matrices de rotation dans le bon ordre, jamais en additionnant des angles
ou des pôles.
\end{solution}

\begin{solution}{exo:b1:plate-kinematics:9}
(a) Dans l'espace des vitesses, placer $a$ à l'origine, $b$ à \qty{30}{mm.an^{-1}} vers
l'est et $c$ à \qty{40}{mm.an^{-1}} vers le sud. (b) $\vec v_{CB} =
\vec v_{CA} - \vec v_{BA} = (0, -40) - (30, 0) = (-30, -40)$ (composantes est,
nord)~: \qty{50}{mm.an^{-1}} vers l'\omterm{def:b1:spherical-astronomy:horizontal}{azimut}
$\arctan\bigl((-30)/(-40)\bigr) + \ang{180} = \ang{216.9}$, sud-ouest. (c) La
dorsale est perpendiculaire à cette direction~: elle est orientée
$\ang{216.9} - \ang{90} = \ang{126.9}$, c'est-à-dire nord-ouest--sud-est.
\end{solution}

\begin{solution}{exo:b1:plate-kinematics:10}
(a) Dans la matrice de Rodrigues, l'identité contribue $3\cos\alpha$, la matrice
$\vec n\,\vec n^{\mathsf T}$ la trace $n_x^2 + n_y^2 + n_z^2 = 1$, fois
$1 - \cos\alpha$, et la matrice antisymétrique rien~: $\operatorname{tr}R =
3\cos\alpha + 1 - \cos\alpha = 1 + 2\cos\alpha$. (b) Ici $\operatorname{tr}R = 0$,
donc $\cos\alpha = -1/2$ et $\alpha = \ang{120}$. L'axe vérifie $R\vec n =
\vec n$~: $R(x, y, z) = (z, x, y)$, donc $x = y = z$ et $\vec n =
(1, 1, 1)/\sqrt3$, à la \omterm{def:b1:spherical-astronomy:geographic}{latitude} $\arcsin(1/\sqrt3) = \ang{35.26}$\,N et à la
\omterm{def:b1:spherical-astronomy:geographic}{longitude} \ang{45}\,E. (c) $R\vec e_x = \vec e_y$, $R\vec e_y = \vec e_z$,
$R\vec e_z = \vec e_x$~: les trois axes sont permutés, comme le fait un tiers de tour autour de
la diagonale. La formule de Rodrigues avec $\alpha = +\ang{120}$ autour de $\vec n$
donne bien $\vec e_x\mapsto -\tfrac12\vec e_x + \tfrac12(1, 1, 1) +
\tfrac{\sqrt3}{2}\cdot\tfrac{1}{\sqrt3}(0, 1, -1) = (0, 1, 0)$, donc la rotation
est antihoraire vue de dessus \ang{35.26}\,N, \ang{45}\,E.
\end{solution}

\begin{solution}{exo:b1:plate-kinematics:11}
(a) Sur une transformante la vitesse relative $\vec\omega\times\vec r$ est le long de la
faille, donc parallèle à $\vec t$~; elle est perpendiculaire à $\vec\omega$, d'où
$\vec t\cdot\vec\omega = 0$. Le plan perpendiculaire à $\vec t$ passant par le
centre contient $\vec r$ (puisque $\vec t$ est horizontal) et $\vec\omega$~: c'est
le plan du \omterm{def:b1:spherical-astronomy:spherical-triangle}{grand cercle} passant par le point perpendiculairement à la faille,
et le pôle est sur ce cercle. Deux tels plans se coupent selon la droite de
$\vec t_1\times\vec t_2$. (b) $\vec t_1 = (0.5436, 0.8367, -0.0673)$ au
\omterm{def:b1:plate-kinematics:triple-junction}{point triple} et $\vec t_2 = (0.4942, 0.8639, -0.0968)$ en Islande, donc
$\vec t_1\times\vec t_2 = (-0.02285, 0.01936, 0.05612)$, de longueur $0.06362$.
Sa \omterm{def:b1:spherical-astronomy:geographic}{latitude} est $\arcsin(0.05612/0.06362) = \ang{61.9}$\,N et sa \omterm{def:b1:spherical-astronomy:geographic}{longitude}
$\ang{180} - \arctan(0.01936/0.02285) = \ang{139.7}$\,E (composante $x$ négative)~:
le pôle du chapitre. (c) Mal~: les deux \omterm{def:b1:spherical-astronomy:spherical-triangle}{grands cercles} se coupent sous un petit
angle, puisque les orientations ne diffèrent que de \ang{8} sur \ang{26} de la limite.
Changer l'une ou l'autre orientation d'un degré déplace le pôle d'environ \ang{15}. Les vraies
déterminations combinent de nombreuses transformantes avec des vitesses d'expansion, dont la variation
le long de la limite fixe la distance angulaire au pôle.
\end{solution}

\begin{solution}{exo:b1:plate-kinematics:12}
(a) $(94 - 4)/3 = \qty{30}{mm.\year^{-1}}$. (b) Les différences à un an d'intervalle sont
$24 - 4 = 20$, $25 - 5 = 20$, $26 - 6 = 20$, $65 - 15 = 50$, $74 - 24 = 50$,
$75 - 25 = 50$, $76 - 26 = 50$, $85 - 65 = 20$ et $94 - 74 = 20$ (mm). Triées~:
cinq valeurs de 20 et quatre de 50~; la médiane est \qty{20}{mm.an^{-1}}. (c) Deux
positions distantes d'un an sont à la même phase de l'oscillation saisonnière,
qui s'annule dans leur différence. Un saut n'affecte que les paires qui l'enjambent,
ici quatre sur neuf~; tant que moins de la moitié des paires enjambent un saut,
la médiane l'ignore, alors que l'estimation par les extrémités en (a) prend tout le saut
pour du mouvement et sort \qty{50}{\percent} trop haut. La méthode échoue quand
les sauts sont si fréquents que la plupart des paires à un an en contiennent un, ou quand la station
bouge de façon irrégulière, par exemple après un grand séisme.
\end{solution}

\begin{solution}{pb:b1:plate-kinematics:1}
\textbf{1.} $\omega = 0.210\times(\pi/180)/\num{e6} = \qty{3.665e-9}{rad.\year^{-1}}$~;
$\omega R = \num{3.665e-9}\times\qty{6.371e9}{mm} = \qty{23.35}{mm.\year^{-1}}$.

\textbf{2.} $\cos\Delta = \sin\ang{61.8}\sin\ang{64.5} + \cos\ang{61.8}
\cos\ang{64.5}\cos(\ang{-157.6}) = 0.79546 - 0.18809 = 0.60737$, donc
$\Delta = \ang{52.60}$.

\textbf{3.} $v = 23.35\sin\ang{52.60} = 23.35\times0.7944 =
\qty{18.55}{mm.\year^{-1}}$.

\textbf{4.} Le numérateur est $\sin\ang{157.6}\cos\ang{61.8} = 0.18011$ et le
dénominateur
\[
  \cos\ang{64.5}\sin\ang{61.8} - \sin\ang{64.5}\cos\ang{61.8}\cos\ang{157.6} =
  0.37940 + 0.39436 = 0.77376;
\]
les deux sont positifs, donc $\beta = \arctan 0.2328 =
\ang{13.1}$. L'Eurasie se déplace vers $\ang{13.1} + \ang{90} = \ang{103.1}$.

\textbf{5.} Est~: $18.55\sin\ang{103.1} = \qty{18.07}{mm.\year^{-1}}$~; nord~:
$18.55\cos\ang{103.1} = \qty{-4.20}{mm.\year^{-1}}$, c'est-à-dire \qty{4.20}{mm.an^{-1}}
vers le sud.

\textbf{6.} $\cos\Delta = 0.8851$, $\Delta = \ang{27.74}$, $v = 23.35
\sin\ang{27.74} = \qty{10.9}{mm.\year^{-1}}$. Les deux plaques tournent autour d'un pôle en
Sibérie orientale~: la dorsale de Gakkel en est proche, et la vitesse diminue comme
$\sin\Delta$ bien que les plaques soient les mêmes.

\textbf{7.} Ouest~: $94/9.779 = \qty{9.61}{mm.\year^{-1}}$~; est~: $97/9.779 =
\qty{9.92}{mm.\year^{-1}}$.

\textbf{8.} $(94 + 97)/9.779 = 191/9.779 = \qty{19.53}{mm.\year^{-1}}$.

\textbf{9.} À partir de C2An~: $54.1/2.581 = \qty{20.96}{mm.\year^{-1}}$~; à partir de la base de
C1n~: $19.8/0.781 = \qty{25.35}{mm.\year^{-1}}$. Elles ne concordent pas~: plus la
limite est jeune, plus la vitesse est élevée. Un vrai changement de vitesse est possible, mais un excès
constant dans les distances, le même à tous les âges, produirait exactement cela.

\textbf{10.} Le rapport des distances à l'ouest et à l'est est de 0.96 à 0.98~: le
flanc est a crû de quelques pour cent de plus. La vitesse totale ne dépend que de la
somme des deux distances, donc l'asymétrie ne la change pas.

\textbf{11.} $(191 - 5.0)/9.779 = \qty{19.02}{mm.\year^{-1}}$~;
$(54.1 - 5.0)/2.581 = \qty{19.02}{mm.\year^{-1}}$~; $(19.8 - 5.0)/0.781 =
\qty{18.95}{mm.\year^{-1}}$. Les trois concordent maintenant à
\qty{0.1}{mm.an^{-1}} près. Les mêmes \qty{5}{km} font un quart des \qty{19.8}{km}
de la paire la plus jeune mais moins de \qty{3}{\percent} des \qty{191}{km} de la
plus ancienne~; les limites anciennes donnent les vitesses les plus fiables.

\textbf{12.} La demi-vitesse corrigée est $186/9.779/2 = \qty{9.51}{mm.\year^{-1}}$,
et $9.51\times0.084 = \qty{0.80}{km}$, un tiers des \qty{2.5}{km} dont
chaque limite est déplacée. On s'attendrait à ce qu'une bande si étroite soit brouillée.
Qu'elle soit reconnue sur la dorsale de Reykjanes, alors qu'elle est absente sur la
dorsale de Gakkel, privée de magma, suggère qu'un apport abondant de
magma enregistre fidèlement le champ même dans des bandes étroites.

\textbf{13.} Les vitesses corrigées, de 18.95 à \qty{19.02}{mm.an^{-1}}, concordent avec
les \qty{19.0}{mm.an^{-1}} du vecteur des plaques~; sans la correction elles le
dépassaient de 3 à 33 pour cent.

\textbf{14.} $\cos\Delta = \sin\ang{64.139}\sin\ang{59.018} +
\cos\ang{64.139}\cos\ang{59.018}\cos\ang{27.554} = 0.97054$, donc $\Delta =
\ang{13.94}$~; la corde vaut $2R\sin(\Delta/2) = 12742\times\sin\ang{6.97} =
\qty{1546}{km}$.

\textbf{15.} $1000/17.5 = 57$ ans.

\textbf{16.} Les deux directions diffèrent de $\ang{123.1} - \ang{114.8} =
\ang{8.3}$~: la composante horizontale le long du \omterm{def:b1:spherical-astronomy:spherical-triangle}{grand cercle} est
$18.9\cos\ang{8.3} = \qty{18.70}{mm.\year^{-1}}$, et le long de la corde
$18.70\cos\ang{6.97} = \qty{18.56}{mm.\year^{-1}}$, soit environ \qty{18.6}{mm.an^{-1}}.

\textbf{17.} La corde s'allonge de \qty{17.5}{mm.an^{-1}}, \qty{6}{\percent}
de moins que les \qty{18.6}{mm.an^{-1}} prédits. Les deux stations se trouvent loin de
la limite de plaques, donc la différence n'est pas locale. Ou bien les plaques bougent
un peu plus lentement aujourd'hui qu'en moyenne sur les trois derniers millions d'années,
ou bien l'une des deux estimations porte une petite erreur systématique~; les mesures
seules ne peuvent trancher. Un \omterm{def:b1:plate-kinematics:euler-pole}{vecteur d'Euler} calculé à partir du seul positionnement par satellite
prédit \qty{17.1}{mm.an^{-1}} pour cette corde, proche de la valeur mesurée.

\textbf{18.} Est~: $21.12 - (-1.07) = \qty{22.19}{mm.\year^{-1}}$~; nord~:
$-6.97 - (-0.97) = \qty{-6.00}{mm.\year^{-1}}$. Vitesse $\sqrt{22.19^2 + 6.00^2} =
\qty{23.0}{mm.\year^{-1}}$, vers $\ang{180} - \arctan(22.19/6.00) =
\ang{105.1}$.

\textbf{19.} La direction est proche des \ang{103} du vecteur des plaques, mais
la vitesse dépasse ses \qty{18.6}{mm.an^{-1}} d'un quart. Les deux stations se trouvent
dans la zone où la limite se déforme, et Höfn, près de la plus grande calotte
glaciaire, bouge avec le rebond de la \omterm{def:b1:layered-earth:layers}{croûte} à mesure que la glace s'amincit~: vers le haut de
\qty{12.9}{mm.an^{-1}}, et latéralement, en s'éloignant de la glace. Des stations
à l'intérieur de la zone de limite mesurent le mouvement des plaques plus une déformation locale.

\textbf{20.} $0.210\times10 = \ang{2.10}$.

\textbf{21.} Le point parcourt un \omterm{def:b1:spherical-astronomy:spherical-triangle}{petit cercle} de rayon $R\sin\Delta$~:
$6371\times0.7944\times(2.10\times\pi/180) = \qty{185.5}{km}$. La largeur
corrigée de \omterm{def:b1:layered-earth:layers}{croûte} formée depuis le bord jeune de C5n, \qty{186}{km} en
\qty{9.8}{\Ma}, est presque la même~: le mouvement a été régulier sur dix
millions d'années.

\textbf{22.} Composantes (est, nord)~: Eurasie
$(22.9\sin\ang{95.1}, 22.9\cos\ang{95.1}) = (22.81, -2.04)$, Nubie
$(19.2\sin\ang{103.8}, 19.2\cos\ang{103.8}) = (18.65, -4.58)$. Eurasie par rapport
à la Nubie~: $(4.16, 2.54)$, c'est-à-dire \qty{4.9}{mm.an^{-1}} vers
$\arctan(4.16/2.54) = \ang{58.6}$, est-nord-est.

\textbf{23.} Un \omterm{def:b1:plate-kinematics:triple-junction}{point triple} de trois dorsales est \omterm{def:b1:plate-kinematics:triple-junction}{stable} quelles que soient les vitesses. La
vitesse de la question 22 fait s'écarter l'Eurasie et la Nubie, d'est-nord-est à
ouest-sud-ouest~: leur limite au \omterm{def:b1:plate-kinematics:triple-junction}{point triple} doit s'ouvrir, lentement~; c'est le
rift de Terceira, dont l'orientation est-sud-est rend l'ouverture oblique.

\textbf{24.} Partie I~: \qty{18.6}{mm.an^{-1}} au centre de l'Islande
(\qty{19.0}{mm.an^{-1}} au profil de la partie II). Partie II~: 18.95 à
\qty{19.02}{mm.an^{-1}} une fois corrigé. Partie III~: la corde Reykjavik--Stavanger
s'allonge de \qty{17.5}{mm.an^{-1}}, là où le vecteur des plaques prédit
\qty{18.6}{mm.an^{-1}}~; la paire à l'intérieur de l'Islande donne \qty{23.0}{mm.an^{-1}},
qui n'est pas le mouvement des plaques. Le vecteur des plaques et les bandes concordent~; les
satellites diffèrent de \qty{6}{\percent}. Les bandes moyennent sur jusqu'à dix
millions d'années, le vecteur des plaques sur \qty{3.16}{\Ma}~; les satellites mesurent
le mouvement des vingt à trente dernières années.

\textbf{25.} Le vecteur des plaques donne \qty{18.55}{mm.an^{-1}} à travers l'Islande, et
les bandes corrigées \qty{19.0}{mm.an^{-1}} à côté~: au millimètre par an près, la \omterm{def:b1:plate-kinematics:spreading-rate}{vitesse totale d'expansion} à travers l'Islande est
\textbf{\qty[text-series-to-math=true]{19}{mm.an^{-1}}}.
\end{solution}
